\section{Reporting Practice: Units of Inference and Sweeping for the Class}
\label{sec:reporting}

Two further findings concern how results are reported rather than what they are, and both cost us
claims.

\paragraph{The unit of inference is equations, not seeds.} A seed-level interval answers ``would a
rerun give this number?'' when the question of interest is ``would other equations give this
number?'' The two differ substantially: per-equation leakage on the 10-equation suite ranges from
$-0.05$ to $1.12$, so resampling equations rather than seeds widens the interval by roughly
$6\times$ (Appendix~\ref{app:units}). We report equation-level bootstraps throughout.

\paragraph{Fix the class, not the instance.} Four times in this work, supplying a non-learned
baseline we had not previously measured reversed a conclusion we had drawn: on a purpose-built
12-law suite, on a 5-law suite, on the $\mathrm{score}_5$ leaderboard (\S\ref{sec:retraction}), and
on the interpretation of a protocol-sensitivity contrast. Each time we had fixed the instance in
front of us without sweeping for the class.

The two suites are worth naming, because we built them specifically to be ``leakage-proof'' and the
audit says they are, at the wrong tier. Both appear in the last block of
Table~\ref{tab:benchmark_survey} at exactly $0\%$ unique variable sets: every law is defined over
the same shared variables, so the variable-identity shortcut we set out to eliminate is genuinely
gone. Both also measure exactly $100\%$ operator-bag separability at tier~2. On the 5-law suite under the joint objective, the consequence is
unambiguous: an \emph{untrained} encoder scores $0.40$ against the trained Tree-LSTM's
$0.32 \pm 0.098$ ($n{=}25$, chance $0.20$); two of the five laws are operator-homogeneous and are
separated under any projection, so every random seed scores exactly $2/5$. We had removed tier~1 and
declared victory without measuring tier~2.

The remedy is mechanical. For \emph{every} protocol under which we state an identification number we
measure all three bag representations and an untrained encoder \emph{under that same protocol}, at
$n{=}25$ seeds (Appendix~\ref{app:sweep}).


The verdicts split cleanly. On the three physics-derived protocols a zero-parameter bag over
variable tokens alone is perfect and the trained encoder is \emph{below} it, so those protocols
cannot license a claim about structure, this paper's thesis applied to its own measurements. Only
the shared-variable design and the two EQNET corpora put the trained encoder above every non-learned
baseline, and those are the only identification claims we advance. Knowing which rows survive is
possible only because the sweep was run for every row, not for the rows we doubted.

One row audits the sweep itself. The Feynman E3 full-token bag reads $0.932$ here against the
$0.992$ of \S\ref{sec:introduction}. These are \emph{not} different protocols: the sweep re-runs
it exactly, and the two disagree only because the held-out paraphrase is redrawn. The estimates are
within one standard deviation, but they straddle the trained encoder's $0.972$, so a comparison
resting on that margin reverses between two runs of the same experiment. Hence the claim in
\S\ref{sec:introduction} is stated against the variable-only bag: $1.000$ with sd $0.000$, perfect
by construction rather than an estimate that happened to land high.
